%Project Summary.

%The Project Summary is a one page, single-spaced summary of the proposed
%activity. The Project Summary should be a self-contained
%description of the proposed research. It should include a statement of
%objectives and methods to be employed. It must address the scientific and/or
%technical merit of the project, for example, the influence that the results
%might have on the direction, progress and thinking in relevant scientific or
%engineering fields.

%It must also address the appropriateness of the proposed methods, and the
%logic and feasibility of the research approach.  The Project Summary should be
%informative to persons working in the same or related fields and, insofar as
%possible, understandable to a scientifically or technically literate lay
%reader.
%%%%%%%%%%%%%%%%%%%%%%%%%%%%%%%%%%%%%%%%%%%%%%%%%%%%%%%%%%%%%%%%%%%%%%%%%%%%%%%%

\begin{singlespace}
  

%The automation of tasks that are dangerous, inconvenient, or otherwise undesirable for humans are one of the main benefits of intelligent and programmable robots. 
Environmental monitoring is a field that could benefit greatly from efforts in automation as the dynamic and spatially transient properties of natural phenomena can be difficult or costly to carry out via traditional human intensive methods. 

We propose research into the design of small environmentally adaptable robotic platforms to expand available methods for environmental monitoring with the objective of increasing the quality and quantity of data collection through greater automation. 

Towards that general goal we identify certain traits to focus on. Multi-domain operation to expand the scope of measurements. Field deployability as platforms are, ideally, designed to function robustly in real-world scenarios. In order to promote use and adoption, these platforms should complement existing methods used by scientists as well as be easy to use, configure, and deploy in the field. 

We propose developing a land-water robot for assisting water quality measurements and oil slick detection to depths of up to 10 meters and capable of expansion via "backpack" modules, such as a water sampling backpack or a volumetric particle counting sensor. 

We also propose the research and development of an amphibious method of locomotion for small robotic platforms. The screwdrive would allow a platform to locomote and transition seamlessly between land and water domains, increasing the robot's adaptibility. This method of locomotion is relatively unexplored, particularly for robotics, yet it is uniquely positioned to fulfill the amphibious space while being simple in design. 

Prior research on large military tank-size screwdrive propulsion systems indicates that optimal screw parameters, such as blade height and pitch, differ for land and water regimes, therefore, characterization at this scale is necessary. Furthermore, we propose the development of a variable-pitch screw to be used in this system both for characterization purposes and for field deployment. One way to achieve this would be the use of origami or kirigami to induce buckling along predefined creases where actuation morphs the geometry of the screw a manner akin to the way an accordion expands and contracts, changing its properties. In order to match the variable screws to their terrain, a terrain classification system is also proposed.  

Finally, noting the benefits of land-water adaptability to the available mission set, we propose development of a land-water-air platform that incorporates principles from the land-water platform, the amphibious powertrain, and adds flight capabilities. This new platform, to the author's knowledge, would be the first of its kind. The addition of flight would allow quicker and easier deployment and retrieval while allowing sensing methodologies not previously available.

%Finally, we propose the development of a land-water-air platform that incorporates principles from the Aquapod, the screwdrive, and adds flight. This new platform, to the author's knowledge, would be the first of its kind. This platform would allow quicker deployment using flight to overcome obstacles while still being able to perform measurement tasks. 

\end{singlespace}